\begin{textsample}{POS Dim 5 – ba – Score 6.00 – ba/ba\_inf\_e\_ef\_2.txt}  \label{ex:f5_pos_274}
APRESENTAÇÃO O Documento Curricular Referencial da Bahia – DCRB para a Educação Infantil e Ensino Fundamental tem como objetivo assegurar os princípios educacionais e os direitos de aprendizagem de todos os estudantes do território estadual, em toda a Educação Básica.
Trata -se de um documento aberto, não prescritivo, que pretende incorporar inovações e atualizações pedagógicas advindas dos marcos legais, do arcabouço teórico-metodológico do currículo, no processo de implementação, considerando, também, aspectos identificados pelos segmentos da comunidade escolar.
Constitui -se numa referência, como o próprio nome deixa antever, para que municípios do Estado da Bahia elaborem os seus currículos com convergência de princípios, intenções e temáticas contidas no Referencial do Estado, para o desenvolvimento de práticas educativas que possibilitem a permanência e o sucesso dos estudantes na escola.
Concretiza -se por meio de sua complementação com os Currículos Escolares e os Planos de Ensino, no âmbito dos Projetos Políticos Pedagógicos ( PPP ) e, também, nas relações entre educadores e estudantes que devem comprometer -se com a aprendizagem como direito do sujeito e dever legal e social de todos.
O Documento Curricular Referencial da Bahia é composto de dois \textbf{volumes} : um para a Educação Infantil e Ensino Fundamental e o outro para o Ensino Médio.
Cabe reforçar que os princípios educacionais nele apresentados são para toda a Educação Básica, independentemente da divisão dos \textbf{volumes}.
Fonte : elaboração Equipe SUPED/SEC – 2019.
PRINCÍPIOS DA EDUCAÇÃO BAIANA DIREITO DE APRENDER PROJETOS POLÍTICO-PEDAGÓGICOS ORGANIZAÇÃO DO TRABALHO PEDAGÓGICO ( PLANOS DE AULA, PROJETOS E \textbf{SEQUÊNCIAS} DIDÁTICAS, ENTRE OUTROS ) DOCUMENTO CURRICULAR REFERENCIAL DA BAHIA PARA A EDUCAÇÃO INFANTIL E ENSINO FUNDAMENTAL Por se tratar de um documento que materializa os “ Atos Curriculares ” que servirão de referência para o Estado, o Documento Curricular Referencial da Bahia para a Educação Infantil e Ensino Fundamental deve assegurar, em todas as suas orientações, sejam elas normativas ou de implementação, os princípios que convergem na educação baiana, expressos nas diretrizes que orientam o Plano Estadual de Educação ( PEE¹ ) : I. erradicação do analfabetismo ; II. universalização do atendimento escolar ; III. superação das desigualdades educacionais, com ênfase no desenvolvimento integral do sujeito, na promoção da cidadania e na erradicação de todas as formas de discriminação ; IV. melhoria da qualidade da educação ; V. formação para o desenvolvimento integral do sujeito, para a cidadania e para o trabalho, com ênfase nos valores morais e éticos nos quais se fundamenta a sociedade ; VI. promoção do princípio da gestão democrática da educação no Estado ; VII. promoção humanística, científica, cultural e tecnológica do Estado ; VIII. valorização dos profissionais da educação ; IX. promoção dos princípios do respeito aos direitos humanos, à diversidade e à sustentabilidade socioambiental.
Com base nesses princípios, o DCRB para a Educação Infantil e Ensino Fundamental dá origem a um processo de contextualização, caracterização e inclusão de especificidades da identidade do Estado da Bahia e seus territórios, considerando uma natureza mais consolidada no próprio documento e, mais particularizada, nos Projetos Políticos Pedagógicos ( PPP ) e Organização do Trabalho Pedagógico ( planos de aula, projetos e \textbf{sequências} didáticas, entre outros ).
Nesse sentido, é fundamental a atuação dos educadores na gestão escolar, no planejamento e na avaliação para que esses fundamentos legais, pactuados pela educação baiana, permeiem o contexto social no qual ocorrem e tornem “ vivas ” as competências e as habilidades que serão apresentadas neste documento.
O DCRB tem como base as orientações normativas da BNCC, complementada à luz das diversidades e das singularidades do território baiano, de modo a colaborar com a (re)escrita dos Projetos Políticos Pedagógicos das Unidades Escolares.
Contém a expectativa de que os estudantes se tornem aptos para enfrentar os desafios contemporâneos, em quaisquer contextos e/ou ambientes em que eles estejam, dentro do território baiano ou fora dele.
Na assunção da BNCC como base, o DCRB destaca as Competências Gerais, descritas no item 2 desse documento, por serem norteadoras da ação educativa para a Educação Básica e suas Modalidades.
Nessa expectativa, é necessário um processo legítimo de valorização da diversidade, conforme evidenciado na competência 9, representada por seus diferentes atores, em todas as Redes de Ensino.
Dessa forma, estudantes, professores, gestores, funcionários, família e comunidade são convidados a desenvolver “ uma escola que faça sentido para o desenvolvimento integral dos estudantes ” nas redes estadual, municipal e privada.
Quanto à sua estrutura, o DCRB apresenta aspectos sobre Territorialidade, Marcos Legais, Marcos Teórico-Conceituais, Marcos Metodológicos, Modalidades da Educação Básica, Temas Integradores e Avaliação Educacional, considerados para toda a Educação Básica, por apresentarem uma fundamentação conceitual e legal que sustenta a política educacional, preparando o “ terreno ” para a compreensão da dimensão sociocognitiva das competências e habilidades, na perspectiva da garantia dos direitos de aprendizagem e desenvolvimento dos estudantes.
O \textbf{Volume} 1 traz as etapas da Educação Infantil e do Ensino Fundamental, apresentando suas etapas e seus \textbf{organizadores} curriculares, por Campo de Experiências e Componentes Curriculares das Áreas do Conhecimento, à luz da Base Nacional Comum Curricular ( BNCC ).
No Organizador Curricular, foram mantidos os códigos alfanuméricos que identificam as aprendizagens, citadas na BNCC, cuja composição é feita por \textbf{letras} e números.
O primeiro par de \textbf{letras} identifica a etapa de ensino, o primeiro par de números identifica o grupo por faixa etária na Educação Infantil e o ano a que se refere à habilidade do Ensino Fundamental.
O segundo par de \textbf{letras} identifica o campo de experiência na Educação Infantil e o Componente Curricular, no caso do Ensino Fundamental.
Por último, mais um par de números que identifica a posição da habilidade na \textbf{numeração} \textbf{sequencial} do campo de experiência na Educação Infantil e, ano ou \textbf{bloco} de ano, no Ensino Fundamental, conforme imagens a seguir : Para identificar habilidades na BNCC que foram contextualizadas para atender às especificidades do território baiano, acrescentou -se o símbolo asterisco ( ) e, para identificar as habilidades elaboradas pela equipe curricular do Estado da Bahia, visando aprendizagem do estudante no reconhecimento da diversidade que o Estado apresenta, foi seguida a estrutura do código alfanumérico da BNCC, acrescentando, ao final, a \textbf{sigla} do Estado ( BA ).
Ainda no Organizador Curricular, destacam -se as competências específicas relacionadas às habilidades que devem ser desenvolvidas pelos estudantes ao longo do Ensino Fundamental.
O \textbf{Volume} 2 trata da etapa do Ensino Médio, assegurando os mesmos princípios e se desdobrando em novos aportes de acordo com as especificidades que a etapa requer.
Também está incutida a organização e Itinerários Formativos para os estudantes nessa fase escolar, o Novo Ensino Médio².
O documento DCRB conclama instituições educacionais e seus educadores a tomá -lo como um conjunto de proposições fundamentadas e justificadas, a serem pensadas com autonomia pedagógica, responsabilidade socioeducacional, atitude formacional criativa e protagonismo institucional, tanto na perspectiva operacional como na vontade de qualificação da formação para a Educação Básica do Estado da Bahia.
A escola tem potencial político, inteligência institucional e operacional para se apropriar das políticas curriculares com responsabilidade e, também, propor singularidades curriculares com pertinência, considerando que o “ direito à aprendizagem e ao desenvolvimento ” dos estudantes e dos professores deverá sempre ser tomado como perspectiva.
Ou seja, pluralmente considerado, por mais que tenha de pleitear como compromisso coletivo “ o bem comum ”, socialmente referenciado, a partir da Educação.
Assim sendo, espera -se que, a partir das práticas e reflexões dos profissionais da Educação, o DCRB possa mobilizar instituições educacionais do Estado, nos seus espaço-temporalidades próprios, para compor renovações diante dos infindáveis, (in)tensos e desafiantes acontecimentos com que as problemáticas humanas e os saberes de possibilidades formativas se apresentam no presente. ¹ Plano Estadual de Educação da Bahia – Lei Estadual nº 13.559, de 11 de maio de 2016. ² Conjunto de políticas, incluídas a BNCC e as DCN, para a ressignificação do Ensino Médio, iniciadas em 2018 pelo Ministério da Educação ( MEC ).

% matched lemmas: bloco, letra, numeração, organizador, sequencial, sequência, sigla, volume
\end{textsample}
